\section{Tracial preliminaries and a free-sum estimate}\label{sec:prelim}
Throughout, $(M,\tau)$ is a von Neumann algebra with faithful normal tracial state. We write
\[
 \normtwo{x}=\tau(x^*x)^{1/2},\qquad \norm{x}=\text{operator norm of }x.
\]
The Hilbert space completion of $M$ in $\normtwo{\cdot}$ is $L^2(M,\tau)$, with inner product $\langle x,y\rangle=\tau(y^*x)$, linear in the first variable. Left multiplication gives its faithful normal representation of $M$. A $*$-isomorphism is a bijective complex-linear map preserving multiplication and adjoints. A linear map between von Neumann algebras is \emph{normal} if it is ultraweakly continuous. We use standard finite von Neumann algebra facts, including bounded Borel functional calculus, Kaplansky density, and the trace-preserving conditional expectation onto a von Neumann subalgebra. The expectation extends to the orthogonal projection of $L^2(M,\tau)$ onto the subalgebra's $L^2$ space. Section~\ref{sec:amplification} recalls the amplification convention separately.

A family $(B_i)_{i\in I}$ of unital subalgebras is \emph{free} if, for every integer $r\geq1$ and every choice of elements satisfying the conditions below,
\[
 \tau(b_1\cdots b_r)=0
 \quad\text{whenever}\quad
 b_j\in B_{i_j},\quad \tau(b_j)=0,\quad
 i_j\ne i_{j+1}\ (1\leq j<r).
\]
Freeness of unital $*$-algebras passes to their generated von Neumann algebras. Indeed, bounded strong approximation gives bounded $L^2$ approximation in a finite tracial algebra; telescoping a product shows continuity of its trace under such approximations. Centering the approximants proves the assertion. Freeness also persists when disjoint collections of the free components are grouped: their algebraic spans consist of scalars and reduced products of centered elements, and the same definition applies to the resulting reduced products.

\subsection{Haar tuples and normal isomorphisms}
\begin{definition}
A \emph{freely generating Haar $d$-tuple} in $M$ is a tuple of unitaries $U=(U_1,\ldots,U_d)$ such that $W^*(U)=M$ and
\[
 \tau(U_g)=0\qquad(g\in\F_d\setminus\{e\}).
\]
Here $U_g$ is the evaluation of the reduced group word $g$ in $U$, with inverse letters evaluated as adjoints.
\end{definition}

\begin{lemma}[Identification by a generating Haar tuple]\label{lem:haar-identification}
For a freely generating Haar $d$-tuple $U$ in $M$, the assignment $\lambda(x_j)\mapsto U_j$ extends to a unital normal trace-preserving $*$-isomorphism $L(\F_d)\to M$.
\end{lemma}
\begin{proof}
The vectors $U_g$, indexed by $g\in\F_d$, are orthonormal because
$\tau(U_h^*U_g)=\tau(U_{h^{-1}g})=\mathbf 1_{g=h}$.
They span $L^2(M,\tau)$. To see this, let $\eta$ be orthogonal to their span. Choose $\eta_m\in M$ converging to $\eta$ in $L^2$. The normal functionals $f_m(x)=\tau(\eta_m^*x)$ converge in norm to $f(x)=\langle x,\eta\rangle$, since $\norm{f_m-f}_{M^*}\leq\normtwo{\eta_m-\eta}$. The space $M_*$ of normal linear functionals is norm closed in $M^*$, so $f$ is normal. It vanishes on the ultraweakly dense word algebra and therefore on all of $M$. Thus $\eta=0$.

Consequently the map $V:\ell^2(\F_d)\to L^2(M,\tau)$ with $V\delta_g=U_g$ is unitary. It intertwines $\lambda(x_j)$ with left multiplication by $U_j$. Conjugation by $V$ identifies their double commutants. The double commutant on the right is the faithful normal left representation of $W^*(U)=M$, so this spatial identification gives the claimed isomorphism. It preserves the trace because $V\delta_e=1$.
\end{proof}

In particular, two freely generating Haar tuples of the same size in $M$ are related by a normal trace-preserving automorphism. The component algebras of a Haar tuple are free: expanding centered Laurent polynomials reduces the assertion to the nonidentity group-word trace, and bounded approximation gives freeness of the von Neumann closures.

For completeness, $L(\F_d)$ is a $\mathrm{II}_1$ factor for $d\geq2$. Its canonical trace is $\tau(x)=\langle x\delta_e,\delta_e\rangle$, which is faithful and normal. Every nonidentity reduced word has infinitely many distinct conjugates: choose a letter $a$ whose powers cause no cancellation at either end of $a^kga^{-k}$. At most two of the $2d$ letters are excluded, and the reduced lengths then tend to infinity. Fourier coefficients of a central element are constant on conjugacy classes and form an $\ell^2$ vector, so every coefficient off the identity vanishes. Faithfulness of the tracial representation implies that the center consists of scalars. The infinitely many orthonormal group unitaries show that the factor is infinite-dimensional.

\subsection{Conjugates of a bounded logarithm}
Fix a freely generating Haar tuple $(A_1,\ldots,A_n,C)$, let $\Gamma=\F_n$, and write $A_g$ for the word in the first $n$ variables indexed by $g\in\Gamma$. Put
\begin{equation}\label{eq:logarithm}
 S=\arg(C),\qquad S_g=A_gSA_g^*\quad(g\in\Gamma),
\end{equation}
where the argument takes values in $(-\pi,\pi]$. The bounded Borel functional calculus for normal operators defines this logarithm in $W^*(C)$; see \cite[Theorem~7.1]{haase2020}. Since the spectral distribution of $C$ is uniform on the circle,
\[
 S=S^*,\quad \norm{S}\leq\pi,\quad \tau(S)=0,\quad e^{\ii S}=C,
 \qquad W^*(S)=W^*(C).
\]

\begin{lemma}[Freeness of the conjugate components]\label{lem:conjugates}
The algebras $A_gW^*(C)A_g^*$, $g\in\Gamma$, form a free family. In particular, each $S_g$ is free from the von Neumann algebra generated by all $S_b$ with $b\ne g$.
\end{lemma}
\begin{proof}
Consider a product of centered elements $A_{g_j}c_jA_{g_j}^*$, with $c_j\in W^*(C)$ and $g_j\ne g_{j+1}$. Cyclicity of the trace rewrites its trace as
\[
 \tau\bigl(c_1 A_{g_1^{-1}g_2}c_2\cdots
          A_{g_{r-1}^{-1}g_r}c_r A_{g_r^{-1}g_1}\bigr).
\]
All interior group elements are nonidentity and hence centered. If the last group element is the identity, omit it; otherwise it too is centered. Freeness of $W^*(A_1,\ldots,A_n)$ and $W^*(C)$ makes the displayed trace zero. For a single factor this is simply $\tau(c_1)=0$. The assertion about the other components follows by grouping free subalgebras.
\end{proof}

For a monomial in $A_j,A_j^*$ and $S$, define its \emph{label} as the product in $\Gamma$ of its group symbols, read in order with each occurrence of $S$ ignored.
\begin{lemma}[Group labels and traces]\label{lem:labels}
A monomial with nonidentity label has trace zero. A monomial with identity label can be expressed as a product of variables $S_g$ (or as the constant $1$).
\end{lemma}
\begin{proof}
For the first assertion, combine consecutive group symbols and consecutive powers of $S$. Remove identity group factors. Split each polynomial in $S$ into its scalar trace and a centered part. A term with all alternating factors centered has trace zero by freeness of the group algebra and $W^*(S)$. Removing a scalar factor permits a reduction to fewer alternating factors, without changing the total group label. Induction on the number of factors ends with either a centered alternating product or a single group element whose label remains nonidentity. Both have trace zero.

For the second assertion, successively use $A_g S=S_g A_g$ to move every group factor to the right. The final group factor is exactly the label, and is therefore $1$.
\end{proof}

\subsection{A bound independent of support size}
For a finitely supported real function $k$ on $\Gamma$, set
\[
 s(k)=\sum_{g\in\Gamma}k(g)S_g.
\]
The coefficient norm $\norm{k}_{\ell^2}$ below is distinguished from the tracial $L^2$ norm $\normtwo{\cdot}$ of operators.
The next estimate uses the creation, annihilation, and diagonal decomposition
underlying free-product Khintchine inequalities; see Ricard and
Xu~\cite[Section~2, Lemmas~2.1--2.3 and Corollary~2.4]{ricard_xu2006}.
We include the scalar self-adjoint argument to give an explicit constant
independent of the number of summands.
\begin{lemma}[Norm of a free sum]\label{lem:free-sum}
For every such $k$,
\begin{equation}\label{eq:free-sum}
 \norm{s(k)}\leq K\norm{k}_{\ell^2},\qquad K=3\pi.
\end{equation}
\end{lemma}
\begin{proof}
We prove the more general bound
\begin{equation}\label{eq:general-free-sum}
 \norm{\sum_j z_j}
 \leq 2\Bigl(\sum_j\normtwo{z_j}^2\Bigr)^{1/2}
       +\max_j\norm{z_j}
\end{equation}
for a finite free family of subalgebras and centered self-adjoint $z_j$ in the respective components.

Work in the tracial $L^2$ space of their generated algebra. Freeness identifies this space with the orthogonal sum of the scalar vector and the reduced centered tensor spaces, whose successive component indices are different. To check the identification, take inner products of reduced centered words. If the two letters at the joining point belong to different components, freeness gives zero. If they belong to the same component, split their product into its trace and its centered part; the latter again has zero contribution, while the scalar term shortens both words. Iteration gives orthogonality for different index sequences and the tensor-product inner product for matching sequences. The reduced word span is dense by the argument used in Lemma~\ref{lem:haar-identification}.

Let $P_j$ project onto the reduced tensors beginning with component $j$, and let $L_{z_j}$ denote left multiplication. On $(1-P_j)L^2$, multiplication prefixes $z_j$, so
\[
 B_j=P_jL_{z_j}(1-P_j),\qquad \norm{B_j}\leq\normtwo{z_j}.
\]
The ranges of the $B_j$ are orthogonal. For every vector $\xi$ this gives
\[
 \norm{\sum_jB_j\xi}^2
   =\sum_j\norm{B_j\xi}^2
   \leq\Bigl(\sum_j\normtwo{z_j}^2\Bigr)\norm{\xi}^2.
\]
The blocks in the opposite direction are $B_j^*$, because $z_j$ is self-adjoint. The block $(1-P_j)L_{z_j}(1-P_j)$ is zero. The remaining diagonal blocks $P_jL_{z_j}P_j$ act on mutually orthogonal subspaces; their sum has norm at most $\max_j\norm{z_j}$. Summing these three parts proves \eqref{eq:general-free-sum}, since the left representation is faithful.

Apply this bound to $z_g=k(g)S_g$. Both $\normtwo{S_g}$ and $\norm{S_g}$ are at most $\pi$, and $\max_g\abs{k(g)}\leq\norm{k}_{\ell^2}$. This proves \eqref{eq:free-sum}.
\end{proof}
